% Part: sets-functions-relations
% Chapter: sets
% Section: proofs-about-sets

\documentclass[../../../include/open-logic-section]{subfiles}

\begin{document}

\olfileid{sfr}{set}{prf}
\olsection{ગણો વિશેની સાબિતીઓ}

\begin{editorial}
આ વિભાગનું સ્થાન \verb|content/methods/proofs|એ લીધું છે;
તેથી મૂળભૂત આયાતક્રમમાં તે આ અધ્યાયમાંથી દૂર કરેલો છે.
\sourcecorrection{OLMD-352}{મૂળની સંપાદકીય નોંધમાં પદ્ધતિઓના ભાગનો માર્ગ એકવચનમાં છે. સ્થિર સ્રોતમાં તે વિભાગ methods/proofs હેઠળ છે; માર્ગ તે મુજબ સુધાર્યો છે.}
\end{editorial}

\begin{explain}
અત્યાર સુધી રજૂ કરેલા ગણો અને તેમના સંકેતો ગણો નિર્ધારિત કરવા
અને તેમની વચ્ચેના સંબંધો વ્યક્ત કરવા અનુકૂળ સંક્ષિપ્ત રીત આપે છે.
આવા સંબંધો અંગે દાવા સાબિત કરવા પણ ઘણી વાર જરૂરી થાય છે.
ગણિતીય સાબિતીઓથી પરિચિત ન હો તો આ નવું લાગશે.
તેથી એક સરળ ઉદાહરણ વિગતે જોઈએ.
કોઈ પણ ગણો $X$ અને $Y$ માટે હંમેશાં
$X \cap (X \cup Y) = X$ થાય તે સાબિત કરીએ.
આવી ગણ-સમાનતા કેવી રીતે સાબિત થાય?
ગણો ફક્ત તેમના !!{element}sથી નક્કી થાય છે:
બે ગણોને એક જ !!{element}s હોય જો અને ફક્ત જો તેઓ સમાન હોય.
એટલે અહીં બતાવવાનું છે કે (ક) $X \cap (X \cup
Y)$નો દરેક !!{element} એ $X$નો પણ !!{element} છે,
અને ઊલટું (ખ) $X$નો દરેક !!{element} એ
$X \cap (X \cup Y)$નો પણ !!{element} છે.
બીજા શબ્દોમાં, બંને (ક) $X \cap (X \cup Y) \subseteq X$
અને (ખ) $X \subseteq X \cap (X \cup Y)$ બતાવીએ.

``$X \cap (X \cup Y)$નો દરેક !!{element}~$z$
એ $X$નો પણ !!{element} છે'' જેવા સામાન્ય દાવાની સાબિતી
માટે મનસ્વી $z \in X \cap (X \cup Y)$ આપેલો છે એમ
પહેલાં ધારીએ, અને તે ધારણાથી $z \in X$ સાબિત કરીએ.
આ રીતને સામાન્ય શરતી સાબિતી તરીકે તમે જાણતા હશો.
ધારણા અને નિષ્કર્ષમાં આવતી વ્યાખ્યાઓ પણ વાપરવી પડશે;
અહીં ખાસ કરીને ``$\cap$'' અને ``$\cup$''ની વ્યાખ્યાઓ.
એટલે (ક) માટેની દલીલ આ પ્રમાણે છે:

\begin{quote}
(ક) પહેલાં $X \cap (X \cup Y) \subseteq X$ બતાવવું છે.
$\subseteq$ની વ્યાખ્યા મુજબ, કોઈ પણ $z$ માટે
જો $z \in X \cap (X \cup Y)$ હોય, તો $z \in X$.
તેથી $z \in X \cap (X \cup Y)$ ધારો.
$z$ બે ગણોના છેદગણનો !!{element} હોય જો અને ફક્ત જો
તે બંને ગણોનો !!{element} હોય; એટલે $z \in X$ અને
$z \in X \cup Y$ મળે છે. ખાસ કરીને $z \in X$,
જે બતાવવું હતું.
\end{quote}

આથી સાબિતીનો પહેલો અડધો ભાગ પૂરો થાય છે.
છેલ્લા પગલામાં એ હકીકત વાપરી છે કે કોઈ ધારણાથી
સંયોજન ($z \in X$ અને $z \in X \cup Y$) ફલિત થાય,
તો દરેક જોડક પણ એ જ ધારણાથી ફલિત થાય.
આ નિયમને સંયોજનનો લોપ, અથવા $\land$નો લોપ, તરીકે જાણતા હશો.
હવે (ખ) સાબિત કરીએ:

\begin{quote}
(ખ) હવે $X \subseteq X \cap (X \cup Y)$ સાબિત કરીએ.
$\subseteq$ની વ્યાખ્યા મુજબ, કોઈ પણ $z$ માટે
જો $z \in X$ હોય, તો $z \in X \cap (X \cup Y)$ પણ થાય.
$z \in X$ ધારો. $z \in X \cap (X \cup Y)$ બતાવવા
``$\cap$''ની વ્યાખ્યા મુજબ (૧) $z \in X$ અને
(૨) $z \in X \cup Y$ બંને બતાવવાં પડશે.
અહીં (૧) ધારણા જ છે, એટલે વધુ કશું સાબિત કરવાનું નથી.
(૨) માટે યાદ રાખો કે $z$ ગણોના યોગગણનો !!{element} હોય
જો અને ફક્ત જો તે તેમાંના ઓછામાં ઓછા એક ગણનો !!{element} હોય.
અહીં $z \in X$ છે અને $X \cup Y$ એ $X$ તથા $Y$નો યોગગણ છે,
એટલે આ શરત મળે છે. તેથી $z \in X \cup Y$.
બંને (૧) $z \in X$ અને (૨) $z \in X \cup Y$ બતાવ્યાં;
એટલે ``$\cap$''ની વ્યાખ્યાથી $z \in X \cap (X \cup Y)$.
\end{quote}

આ દલીલ લાંબી હતી, પણ ગણો અને તેમના સંબંધો વિશે કેવી રીતે
વિચારીએ તે બતાવે છે. સામાન્ય રીતે આટલી વિગત આપતા નથી;
ખાસ કરીને બધી વ્યાખ્યાઓ ફરી કહેતા નથી.
વધુ આગળના પાઠમાં આ પરિણામની સાબિતી ઘણી સંક્ષિપ્ત હોય.
તે આ પ્રમાણે હોઈ શકે.
\end{explain}

\begin{prop}[અવશોષણ]
કોઈ પણ ગણો $X$, $Y$ માટે,
\[
X \cap (X \cup Y) = X
\]
\end{prop}

\begin{proof}
(ક) ધારો કે $z \in X \cap (X \cup Y)$.
તો $z \in X$, એટલે $X \cap (X \cup Y) \subseteq X$.

(ખ) હવે $z \in X$ ધારો.
તો $z \in X \cup Y$ પણ થાય અને તેથી
$z \in X \cap (X \cup Y)$ પણ થાય.
આમ $X \subseteq X \cap (X \cup Y)$.
\end{proof}

\begin{prob}
$X \cup (X \cap Y) = X$ વિગતે સાબિત કરો.
પછી તેની ટૂંકી, સંક્ષિપ્ત સાબિતી આપો.
(સૂચન: $X \cup (X \cap Y) \subseteq X$ની દિશા માટે
કિસ્સાવાર સાબિતી, એટલે કે $\lor$નો લોપ, જોઈએ.)
\end{prob}

\end{document}

